\originalsection{作业 12: 滚动与冲量}
\begin{exercise}
原题 1. 半径 $R$ 自行车轮在水平面无滑向右匀速 $V$. 轮缘珠初始在最高点 $x=x_0,y=2R$. 求地面系珠位置分量.
\end{exercise}
\begin{solution}
轮心 $x_C=x_0+Vt,y_C=R$, 顺时针角 $\phi=Vt/R$, 顶部珠相对中心分量为 $(R\sin\phi,R\cos\phi)$. 故 $x=x_0+Vt+R\sin(Vt/R)$, $y=R+R\cos(Vt/R)$.
\end{solution}
\begin{exercise}
原题 2. 半径 $b$、质量 $M$ 球在半径 $R>b$ 的环内侧无滑滚动. 图示球自转为顺时针角速率 $\omega$, 环逆时针为 $\Omega$, 球心沿环方向运动. 求球心速率.
\end{exercise}
\begin{solution}
取环逆时针切向为正. 接触处球的速度为 $v_C-b\omega$, 环面接触速度为 $R\Omega$. 两物质点无滑要求相等, 所以 $v_C=R\Omega+b\omega$. 球心轨道半径 $R-b$ 只决定球心轨道角速度 $v_C/(R-b)$, 不替换环面线速度中的 $R$.
\end{solution}
\begin{exercise}
原题 3. 均匀立方体质量 $m$、边长 $s$, 在光滑桌上以速率 $v_0$ 水平滑行, 与低凸缘瞬时非弹性碰撞后绕凸缘无摩擦转动. 质心轴惯量 $ms^2/6$, 重力向下. (a) 碰后初角速度. (b) 不翻过凸缘的最大初速.
\end{exercise}
\begin{solution}
凸缘到质心距离 $s/\sqrt2$, $I_P=ms^2/6+ms^2/2=2ms^2/3$. 碰前角动量 $mv_0s/2$, 所以 (a) $\omega_a=3v_0/(4s)$. 质心越过支点正上方需上升 $s(1/\sqrt2-1/2)$, 碰后动能 $3mv_0^2/16$. (b) 阈值 $v_{0,\max}=\sqrt{(8gs/3)(\sqrt2-1)}$, 严格低于它会倒回, 等号到达不稳定最高点的极限情形.
\end{solution}
\begin{exercise}
原题 4. 悠悠球质量 $m$, 外半径 $R$, 轴半径 $R/3$, $I_0=mR^2/2$, 在倾角 $\phi$ 的坡面上. 绳从小轴下侧沿坡向上牵拉 $F$, 静摩擦系数 $\mu_s$. 求无滑滚动允许的最大 $F$.
\end{exercise}
\begin{solution}
取沿坡向上为正, 摩擦分量 $f$ 可正可负. $F+f-mg\sin\phi=ma$. 以出页逆时针为角正, 下侧出绳产生正矩 $FR/3$, 摩擦矩 $Rf$, 无滑向上滚动 $\alpha=-a/R$. 因而 $FR/3+Rf=-(mR/2)a$. 消元得 $a=4F/(9m)-2g\sin\phi/3$, $f=mg\sin\phi/3-5F/9$. 由 $|f|\le\mu_smg\cos\phi$,
\begin{equation}
F_{\max}=\frac35mg(\sin\phi+3\mu_s\cos\phi).
\end{equation}
还存在下界 $F\ge(3/5)mg(\sin\phi-3\mu_s\cos\phi)$, 若下界为负则取 $F\ge0$. 仅给最大值不能暗示全部更小牵引力一定无滑.
\end{solution}
\begin{exercise}
原题 5. 均匀台球 $m,R,I_C=2mR^2/5$, 初始静止. 球杆在高于球心 $h$ 的位置给水平强冲量, 碰撞期间摩擦可忽略, 球得到质心速率 $v_0$ 和角速度 $\omega_0$. 此后动摩擦 $\mu_k$ 使球滑转, 达到无滑时最大速率为 $9v_0/7$. 求 $h/R$.
\end{exercise}
\begin{solution}
冲量 $J=mv_0$, 顺时针角速度大小 $|\omega_0|=Jh/I_C=5v_0h/(2R^2)$. 无滑最终关系可由摩擦平动与转动方程消去时间: $v_f=[v_0+(2/5)R|\omega_0|]/(1+2/5)=5v_0(1+h/R)/7$. 给定 $v_f=9v_0/7$ 得 $h/R=4/5$. 此时初始接触点向后滑, 摩擦向前, 所以质心确实加速而自转减速.
\end{solution}
